\begin{enumerate}
	\item \textbf{Affirmation fausse.}

	Dans une suite arithmétique de raison $r$, on passe d'un terme au suivant en
	ajoutant $r$ :
	\[u_6 = u_5 + 5 = 15 \qquad\text{et}\qquad u_7 = u_6 + 5 = 20.\]

	\item \textbf{Affirmation fausse.}

	On calcule les premiers termes : $u_0 = 1$, $u_1 = 3$, $u_2 = 9$, $u_3 = 19$.
	\[\dfrac{u_1}{u_0} = 3, \qquad \dfrac{u_2}{u_1} = 3 \qquad
	\text{mais} \qquad \dfrac{u_3}{u_2} = \dfrac{19}{9} \neq 3.\]
	Le quotient de deux termes consécutifs n'est pas constant : la suite $(u_n)$ n'est pas géométrique.

	\item \textbf{Affirmation vraie.}

	L'équation de la tangente au point d'abscisse $a$ est $y = f'(a)(x - a) + f(a)$.

	$f(1) = 5 - 2 + 3 = 6$ ; $f'(x) = 10x - 2$, donc $f'(1) = 8$. La tangente a pour équation
	\[y = 8(x - 1) + 6 \]
	\[ y = 8x - 8 + 6 \]
	\[ y = 8x - 2\]

	\item \textbf{Affirmation fausse.}

	Une diminution de 15\,\% correspond au coefficient multiplicateur
	$1 - \dfrac{15}{100} = 0,85$.

	Chaque année, il reste 85\,\% des truites de l'année précédente, auxquelles on ajoute les
	100 truites réintroduites : pour tout entier naturel $n$, $u_{n+1} = 0,85u_n + 100$.
\end{enumerate}
